\section{Proofs for Coarse Channels}
\label{sec:proofs-coarse-channels}

\noindent\emph{Proof of \Cref{prop:cg-finest-independent-partition}.}
Suppose $P$ factors over partitions $\mathfrak p$ and $\mathfrak q$. Marginalizing
the second factorization to a block $I\in\mathfrak p$ gives
\begin{equation*}
P_I=\bigotimes_{J\in\mathfrak q:I\cap J\ne\varnothing}P_{I\cap J}.
\end{equation*}
Substitution in $P=\bigotimes_{I\in\mathfrak p}P_I$ proves factorization over the
common refinement. The finite set of independent partitions is nonempty, since it
contains $\{V\}$, and is closed under common refinement. Its common finest refinement
is independent and uniquely refines every member. Every independent partition has
negative objective in \eqref{eq:cg-independent-selector-objective}, while every
dependent partition has positive objective. Among the former the finest has the
unique largest block count. If it is intermediate, refinement from $\{V\}$ toward it
strictly decreases the objective; any further strict refinement is dependent and
raises the objective across zero. Relabeling permutes the marginal products, and
separate bimeasurable changes push each factor to its corresponding marginal, so
independence and its finest partition have the asserted naturality and invariance.

For the conditional assertion there are finitely many partitions and pairs of
partitions. For each one with $\mathsf T_{P\mid\Upsilon}=0$, its conditional product
identity holds on a set of full $P_\Upsilon$-mass. Intersect these finitely many sets.
On that intersection the preceding marginalization argument proves closure under
common refinement, hence the unique finest conditional partition and its objective
characterization. The conditional laws given $\Upsilon=Y_V$ are point masses and
therefore products over every partition. These constructions choose a factorization
objective; no identity with the evidence bound has been introduced. $\square$

\noindent\emph{Proof of \Cref{prop:cg-independent-selector-instability}.}
The exponent in \eqref{eq:cg-independent-pair-law} is a sum of terms on disjoint
pairs, so its normalized law is the product of its pair marginals. Within a pair
the odds ratio is $e^{4\mathrm j}\ne1$, so the pair is not independent and cannot
split in an independent partition. This gives its finest partition. For
\eqref{eq:cg-independent-pair-perturbation} all edges of the path have positive
coupling. Any nontrivial partition cuts an edge $ij$. Conditional on all other
coordinates, the two-spin odds ratio across this edge is $e^{4\mathrm j_{ij}}\ne1$,
where $\mathrm j_{ij}$ is $\mathrm j$ or $\mathrm r$. A product over the partition
would instead make the two spins conditionally independent, with ratio one. Positivity
of the finite law permits conditioning on every outside assignment. Thus no nontrivial
independent partition exists. Its finitely many normalized atom probabilities are
continuous in $\mathrm r$ down to zero, which gives total-variation convergence. $\square$

\noindent\emph{Proof of \Cref{prop:cg-exchangeable-hard-selector}.}
Because every agent permutation leaves the input law unchanged, naturality forces the
output unlabeled partition to be invariant under every permutation. If it contained
a block $I$ with $2\leq|I|\leq N-1$, choose distinct $i,j\in I$ and $k\notin I$.
The transposition of $j,k$ fixes $i$ and sends $I$ to a different block containing $i$.
Two distinct blocks of a partition cannot intersect. Therefore no such block exists,
leaving only the singleton and one-block partitions. The argument concerns a
deterministic partition from this law alone and imposes no restriction on a
permutation-invariant probability distribution over partitions. $\square$

\noindent\emph{Proof of \Cref{prop:cg-partition-information-stability}.}
For a partition $\mathfrak p$ of $I$, the entropy numerator satisfies
\begin{equation*}
\sum_{J\in\mathfrak p}\mathrm H(P_J)-\mathrm H(P_I)
=\KL\Bigl(P_I\Bigm\Vert\bigotimes_{J\in\mathfrak p}P_J\Bigr)
\end{equation*}
by \cref{thm:elbo-total-correlation-chain}. It is nonnegative and is zero exactly at a
product. The minimum ranges over a finite set, so its zero case is a nontrivial
independent partition.

For each $J\subseteq I$, marginalization gives
$\operatorname{TV}(P_J,Q_J)\leq\delta$. Couple these two finite-alphabet variables
with mismatch probability $\delta_J\leq\delta$. Writing the mismatch indicator as
an auxiliary binary coordinate and applying the entropy chain rule gives
\begin{equation*}
|\mathrm H(P_J)-\mathrm H(Q_J)|
\leq h_2(\delta_J)+\delta_J\log D_J
\leq h_2(\delta)+\delta\log D_J.
\end{equation*}
The last step uses $\delta\leq1/2$. A singleton alphabet contributes zero. For a
partition with $k\geq2$ blocks, $\sum_{J\in\mathfrak p}\log D_J=\log D_I$;
therefore its normalized entropy expression changes by at most
\begin{equation*}
\frac{(k+1)h_2(\delta)+2\delta\log D_I}{k-1}
\leq3h_2(\delta)+2\delta\log D_I.
\end{equation*}
The minima of finitely many paired expressions differ by at most the largest
paired difference. This proves \eqref{eq:cg-partition-information-stability}.
The bound with $D_V\geq D_I$ proves the two threshold inclusions. A score at distance
greater than $\eta$ from the threshold cannot cross it, so the entire qualified
subset family and its maximal members are unchanged. Finally, the relative-entropy
representation of each numerator is invariant under separate bimeasurable changes
by \cref{thm:cg-dpi-equality} applied in both directions. The finite-alphabet entropy
bound supplies no corresponding bound on unrestricted continuous laws. $\square$

The proofs for \cref{sec:cg-retained-label-asymptotics} use the notation fixed at the head of
that section, including its convention on the constant $\mathrm C_{\mathsf K}$.

\noindent\emph{Proof of \Cref{prop:cg-descent-common-recovery}.}
For a law $P\ll P_0$ with density $r$ and bounded measurable $g$ on $\mathsf Y$,
$\int g d(PK)=\E_{\mathbb J_{P_0}}[g(Y)r(X)]=\E_{\mathbb J_{P_0}}[g(Y)\bar r(Y)]$ with
$\bar r(Y)=\E_{\mathbb J_{P_0}}[r(X)\given Y]$, because $r$ is $P_0$-integrable. This is the
computation in the proof of \Cref{thm:cg-dpi-equality}, which uses no finiteness of relative entropy,
and it shows that a measurable function of the coarse coordinate that agrees with $r$ almost everywhere
is $\bar r=d(PK)/d(P_0K)$. (i) If $r$ descends, disintegrating $\mathbb J_{P_0}$ over $y$ gives
$(PKR_{P_0})(A)=\int\bar r(y)R_{P_0}(y,A)(P_0K)(dy)=\E_{\mathbb J_{P_0}}[\mathbf 1_A(X)\bar r(Y)]
=\E_{\mathbb J_{P_0}}[\mathbf 1_A(X)r(X)]=P(A)$, and $r=1$ gives $P_0KR_{P_0}=P_0$. Conversely, let one
Markov kernel $R$ recover $P_0$ and $P$. Conditional Jensen for the concave square root gives
$\E_{\mathbb J_{P_0}}[\sqrt{r(X)}\given Y]\leq\sqrt{\bar r(Y)}$ almost surely, so the affinity
$\E_{P_0}[\sqrt r]$, which is at most $1$ by the Cauchy-Schwarz inequality, does not decrease through
$K$: $\E_{P_0}[\sqrt r]\leq\E_{P_0K}[\sqrt{\bar r}]$. The kernel $R$ sends the pair $(PK,P_0K)$, of density
$\bar r$, to $(P,P_0)$, of density $r$, so the same inequality through $R$ gives
$\E_{P_0K}[\sqrt{\bar r}]\leq\E_{P_0}[\sqrt r]$. The two finite affinities are therefore equal, the
conditional inequality is an almost sure equality, and
\begin{equation*}
\E_{\mathbb J_{P_0}}\bigl[(\sqrt{r(X)}-\sqrt{\bar r(Y)})^2\bigr]
=2-2\E_{\mathbb J_{P_0}}\bigl[\sqrt{\bar r(Y)}\E_{\mathbb J_{P_0}}[\sqrt{r(X)}\given Y]\bigr]
=2-2\E_{\mathbb J_{P_0}}[\bar r(Y)]=0,
\end{equation*}
where $\sqrt{r(X)\bar r(Y)}$ is integrable because it is at most $(r(X)+\bar r(Y))/2$. Hence $r$
descends. Applied to each $P_\theta$, and to a single law $P\ll P_0$, this proves the first sentence of
(i) and the characterization of the laws that $R_{P_0}$ recovers. The functions in $L^1(P_0)$ that
agree $\mathbb J_{P_0}$-almost everywhere with a measurable function of the coarse coordinate form a
linear subspace containing $1$ and, when every $r_\theta$ descends, every $r_\theta$; it is closed, so
it then contains their closed linear span: if $r^{(n)}\to r$ in $L^1(P_0)$ and $r^{(n)}(X)=\bar r^{(n)}(Y)$ almost surely, then, since $X$
has law $P_0$ under
$\mathbb J_{P_0}$, a subsequence converges $\mathbb J_{P_0}$-almost surely, and the limit superior of
$\bar r^{(n)}$ along it, set to zero where it is infinite, is a measurable function of $y$ equal to
$r(x)$ almost everywhere. For records, $d\bar{\mathsf m}_o/d\varrho=L_o/M_o$ by
\Cref{def:rg-measure-pair}, and $d(\bar{\mathsf m}_oC)/d(\varrho C)=L^c_o/M_o$ by
\eqref{eq:rg-reference-evidence-pair}. (ii) The graph $\operatorname{gr}c=\{(x,y):c(x)=y\}$ is
measurable because the diagonal of the standard Borel space $\mathsf Y$ is, so
$y\mapsto R(y,c^{-1}(\{y\}))=\int\mathbf 1_{\operatorname{gr}c}(x,y)R(y,dx)$ is measurable. Under $K_c$
the joint $\mathbb J_{P_0}$ is the law of $(X,c(X))$ and is carried by $\operatorname{gr}c$, which gives
the stated form of descent, and $\int R_{P_0}(y,c^{-1}(\{y\}))(c_\#P_0)(dy)=\mathbb J_{P_0}(\operatorname{gr}c)=1$,
so every version of $R_{P_0}$ is supported on fibers and recovers $P_0$. If $R$ is supported on fibers
and recovers $P_0$, the measure on $\mathsf X\times\mathsf Y$ that gives $A\times B$ the mass
$\int_BR(y,A)(c_\#P_0)(dy)$ is carried by $\operatorname{gr}c$ and has first marginal $P_0K_cR=P_0$, so
its value on $A\times B$ equals its value on $(A\cap c^{-1}(B))\times\mathsf Y$, which is
$P_0(A\cap c^{-1}(B))=\mathbb J_{P_0}(A\times B)$. Hence $R$ disintegrates $\mathbb J_{P_0}$ over $y$ and
is a version of $R_{P_0}$, and two versions agree $c_\#P_0$-almost everywhere because $\mathsf X$ is
countably generated. Since $P_\theta K_c\ll P_0K_c$, all versions give the same $P_\theta K_cR$, and (i)
gives the rest. (iii) The fiber of $c\times\operatorname{id}_{\mathsf X_2}$ over $(y,x_2)$ is
$c^{-1}(\{y\})\times\{x_2\}$, and the first marginal of $(c\times\operatorname{id}_{\mathsf X_2})_\#P_0$
is the law of $c(X_1)$, so $R$ is supported on fibers and, by (ii), recovers $P_0$ and every $P_\theta$
exactly when it recovers $P_0$ and every $r_\theta$ descends. For measurable $A_1\subseteq\mathsf X_1$,
$A_2\subseteq\mathsf X_2$, and $B\subseteq\mathsf Y$, fiber support gives
$R'(c(X_1),A_1\cap c^{-1}(B))=R'(c(X_1),A_1)\mathbf 1_B(c(X_1))$ almost surely, and measurable
rectangles generate the product $\sigma$-algebra, so $R$ recovers $P_0$ exactly when
\begin{equation*}
P_0\bigl(X_1\in A_1,c(X_1)\in B,X_2\in A_2\bigr)
=\E_{P_0}\bigl[R'(c(X_1),A_1)\mathbf 1_B(c(X_1))\mathbf 1_{A_2}(X_2)\bigr]
\end{equation*}
for all such sets. With $A_2=\mathsf X_2$ this says that $R'$ is a regular conditional law of $X_1$ given
$c(X_1)$. Granted that, conditioning both sides on $c(X_1)$ shows that the identity is equivalent to
$P_0(X_1\in A_1,X_2\in A_2\given c(X_1))=P_0(X_1\in A_1\given c(X_1))P_0(X_2\in A_2\given c(X_1))$
almost surely for all $A_1,A_2$, which is the conditional independence. These conditions on $P_0$
and the $r_\theta$ are equivalent to adequacy in the sense named after the proposition, that is, to a
conditional law of $X_1$ given $c(X_1)$ common to $P_0$ and every $P_\theta$ together with the
conditional independence of $X_1$ and $X_2$ given $c(X_1)$ under each of these laws. Conditional
independence under $P_0$ with descent $r_\theta(x_1,x_2)=\bar r_\theta(c(x_1),x_2)$ makes the
conditional law of $X_1$ given $(c(X_1),X_2)$ under $P_\theta$, whose density with respect to $P_0$ is a
function of $(c(X_1),X_2)$, equal to that under $P_0$, which is the conditional law of $X_1$ given
$c(X_1)$ under $P_0$, so the conditional independence holds under every $P_\theta$ with this common
conditional law; conversely, adequacy writes each member of the family as the law of
$(c(X_1),X_2)$ followed by the common kernel to $X_1$, so every $r_\theta$ is the density between
these laws of $(c(X_1),X_2)$ evaluated at $(c(x_1),x_2)$ and descends. In the witness every $R'$ is
supported on the one fiber, the density $1$ of the family descends,
$R_{P_0}((y,x_2),\cdot)=\delta_{(x_2,x_2)}$ recovers $P_0$,
and $R'(y,\cdot)\otimes\delta_{x_2}$ sends $P_0$ to the product of $R'(y,\cdot)$ with the uniform law on
$\{0,1\}$, which gives $\{x_1\neq x_2\}$ the mass $\tfrac12R'(y,\{1\})+\tfrac12R'(y,\{0\})=\tfrac12$.
$\square$

\noindent\emph{Proof of \Cref{thm:rg-exact-coarse-vfe}.}
(i) Since $C(y,\mathsf Z)=1$, a bounded function of $o$ alone has the same integral against
$\mathbb P^c$ and against $\mathbb P$, which gives normalization and the observation marginal. For
bounded measurable $f$ on $\mathsf Z$ and $g$ on $\mathsf O$ put $h(y)=\int f(z)C(y,dz)$. Then
\begin{align*}
\int g\Bigl(\int f d(\boldsymbol\Pi_oC)\Bigr)\mathbb P^c(do\times\mathsf Z)
&=\int g\Bigl(\int h d\boldsymbol\Pi_o\Bigr)\mathbb P(do\times\mathsf Y)\\
&=\int g(o)h(y)\mathbb P(do,dy)
=\int g(o)f(z)\mathbb P^c(do,dz),
\end{align*}
using the preserved marginal, the defining property of the fine version, and
\eqref{eq:rg-common-pushforward}. Hence $o\mapsto\boldsymbol\Pi_oC$ is a posterior version of
$\mathbb P^c$ at every observation where it is declared, because $C$ acts on the conditioned
variable only. If $\boldsymbol\Pi^c_o(N)=0$, then $C(\cdot,N)=0$ $\boldsymbol\Pi_o$-almost surely,
hence $\mathbb Q_o$-almost surely under $\mathbb Q_o\ll\boldsymbol\Pi_o$, and $\mathbb Q^c_o(N)=0$.
(ii) Disintegrated over $y$, both lifts have the marginals $\mathbb Q_o$ and $\boldsymbol\Pi_o$ and
the common conditional $C(y,\cdot)$, so the chain rule of \Cref{lem:prob-kl-chain-rule}, in
$[0,+\infty]$, gives $\KL(\widehat{\mathbb Q}_o\Vert\widehat{\boldsymbol\Pi}_o)
=\KL(\mathbb Q_o\Vert\boldsymbol\Pi_o)$; disintegrated over $z$, it gives
$\KL(\widehat{\mathbb Q}_o\Vert\widehat{\boldsymbol\Pi}_o)
=\KL(\mathbb Q^c_o\Vert\boldsymbol\Pi^c_o)+\Delta(o)$. Adding the finite real $-\log p(o)$, the
same at both scales by (i), and using that $\boldsymbol\Pi^c_o$ is the coarse posterior version gives
\eqref{eq:rg-vfe-chain-rule}; every term lies in $[0,+\infty]$ before the finite shift, so no
infinite quantity is subtracted. (iii) $\Delta(o)$ is the integral of a nonnegative measurable
function, which vanishes exactly when the integrand vanishes $\mathbb Q^c_o$-almost surely, and a
relative entropy vanishes exactly when its arguments agree. (iv) Finite fine divergence makes both
terms on the right of \eqref{eq:rg-vfe-chain-rule} finite. If $\Delta(o)=0$, the two divergences
are equal and finite, and \Cref{thm:cg-dpi-equality,cor:cg-pairwise-bayes-recovery} give recovery
of both laws by the Bayes kernel of $\boldsymbol\Pi_o$ under $C$. Conversely, a kernel $R$
recovering both laws gives $\KL(\mathbb Q_o\Vert\boldsymbol\Pi_o)\geq\KL(\mathbb Q^c_o\Vert\boldsymbol\Pi^c_o)
\geq\KL(\mathbb Q^c_oR\Vert\boldsymbol\Pi^c_oR)=\KL(\mathbb Q_o\Vert\boldsymbol\Pi_o)$ by data
processing through $C$ and then $R$, so $\Delta(o)=0$. $\square$

\noindent\emph{Proof of \Cref{thm:rg-pointwise-parent-datum}.}
(i) is \Cref{thm:rg-exact-coarse-vfe} with $\mathbb P=\mathbb P_I$, $C=C_A$, and the admitted
recognition law. (ii) On standard Borel spaces a disintegration of $\mathbb P_A(\cdot\given X)$
over $(m_A,\xi_A)$ exists and is jointly measurable (\Cref{ch:probability}); it satisfies the
compatibility identity by construction, while a family declared
independently of $\mathbb P_A$ satisfies it only if assumed. (iii) The pushforward formulas hold by
testing each parent law with functions of one coordinate. For two non-Dirac marginals,
\Cref{prop:prob-marginals-do-not-determine-joint} gives distinct couplings with the same pair, and
tensoring each with one fixed law on $\boldsymbol\Xi_A\times\mathsf H_A$ gives distinct laws on
$\mathsf Z_A$. For two points $h\neq h'$ of $\boldsymbol\Xi_A\times\mathsf H_A$ and the
$(b_A,m_A)$ marginal $R$ of the given law, $R\otimes\delta_h$ and $R\otimes\delta_{h'}$ share the
pair and differ, so one of them differs from the given law. If that space is a point and one
marginal is Dirac, every coupling assigns mass one to the Dirac coordinate. $\square$

\noindent\emph{Proof of \Cref{prop:cg-equivariant-bayes-versions}.}
(i) By equivariance,
$\mathbb J_{g_\#Q}(A\times B)=\mathbb J_Q(g^{-1}A\times g^{-1}B)$. Bimeasurability makes $R_Q^g$ a
Markov kernel, the $\mathsf Y$-marginal of $\mathbb J_{g_\#Q}$ is $g_\#(QK)$ by
\Cref{prop:cg-equivariant-channels}, and the change of variables $y=gy'$ gives
$\int\mathbf 1_B(y)R_Q^g(y,A)(g_\#(QK))(dy)=\mathbb J_{g_\#Q}(A\times B)$. (ii) With
$(g_\#P)K=g_\#(PK)$, $((g_\#P)KR_Q^g)(A)=(PKR_Q)(g^{-1}A)=(g_\#P)(A)$; the relative-entropy
identities are invariance under bimeasurable bijections, \eqref{eq:cg-rn-bimeasurable} and
\eqref{eq:cg-kl-bimeasurable-invariance}. (iii) Where $(QK)(\{y\})>0$ the first value is the value
of every version; the remaining points are $QK$-null, both branches transform as stated because
$(g_\#Q)(\{gx\})=Q(\{x\})$ and $K(gx,\{gy\})=K(x,\{y\})$, and the first branch is a rational
function of $Q$ on a Borel set. (iv) Each $R_Q^h$ with $h\in\mathbb G_Q$ is a version of $R_Q$ by
(i). Joint measurability makes $(h,y)\mapsto R_Q^h(y,A)$ measurable, so $\bar R_Q$ is a Markov
kernel by Tonelli's theorem and monotone convergence, and the defining identity, linear in the
kernel, passes to the mixture. Left invariance of $dh$ gives $\bar R_Q(ky,kA)=\bar R_Q(y,A)$ for
$k\in\mathbb G_Q$. If $g_\#Q=g'_\#Q$ then $g^{-1}g'\in\mathbb G_Q$, so $\bar R_Q^{g'}=\bar R_Q^g$;
finally $\bar R_Q^{kg}(ky,kA)=\bar R_Q^g(y,A)$. $\square$

\noindent\emph{Proof of \Cref{prop:rg-scale-functoriality}.}
(i) Functoriality forces the stated composites, and associativity gives
$C_{kr}\circ C_{\ell k}=C_{\ell r}$. (ii) is the functor axiom on two routes, and for a rule
family it is the axiom on the arrows indexed by $b_1$, $b_2$, and $b_1b_2$. (iii) Each composite of
configuration-keeping or additive arrows is the same statistic of the final partition, and the
witness is direct evaluation: by pairs $(0,0,1,0,1,0)\mapsto(0,1,1)\mapsto1$, by triples
$(0,0,1,0,1,0)\mapsto(0,0)\mapsto0$. (iv) If stochastic matrices $P,Q\geq0$ are mutually inverse,
rank gives them equal sizes. From $PQ=I$, every positive entry $P_{ik}$ forces row $k$ of $Q$ to
vanish off column $i$, so each column of $P$ has at most one positive entry, and unit row sums make
$P$ a permutation matrix. $\square$

\noindent\emph{Proof of \Cref{prop:rg-thermo-finite-labeled-fisher}.}
Differentiating $\sum_xp_\gamma(x)=1$ on the fixed support gives the branch-score identity, and
differentiating the softmax gives $\nabla\log W_\gamma=\xi_\gamma-\bar\xi$, hence
\eqref{eq:rg-thermo-finite-complete-score}. In the outer product of
\eqref{eq:rg-thermo-finite-complete-score}, summation over $x$ within each branch removes both cross
terms by the branch-score identity, which gives \eqref{eq:rg-thermo-finite-labeled-fisher}; the
label marginal has score $\xi_\gamma-\bar\xi$, so its Fisher tensor is
$\operatorname{Cov}_W(\xi_\gamma)$. Positive denominators on a finite support make the $C^3$ masses
give a $C^2$ Fisher tensor, and the Leibniz rule gives
\eqref{eq:rg-thermo-finite-weighted-fisher-derivatives}. A common shift cancels from
the weights and adds one vector to every $\xi_\gamma$, which leaves the covariance unchanged as a
function of $\theta$. Finally $v^\top g^F(P)v=\sum_zP(z)(v^\top\mathbf s_P(z))^2\geq0$. $\square$

\noindent\emph{Proof of \Cref{prop:rg-thermo-finite-missing-information}.}
On its positive support $\mathbf s^c=\mathbf s_P+\nabla\log K^{\mathrm{obs}}(o\given\gamma,x)$, and
each record row has fixed support and sums to one, so
$\sum_oK^{\mathrm{obs}}(o\given\gamma,x)\nabla\log K^{\mathrm{obs}}(o\given\gamma,x)=0$ and the cross
terms of the second moment vanish; this is
\eqref{eq:rg-thermo-finite-complete-observation-fisher}. The record is the image of the complete
sample under the parameter-independent map $(\gamma,x,o)\mapsto o$, with conditional law
$\boldsymbol\Pi_o$ given $o\in\mathsf O^+$. A finite fixed-support family with $C^3$ masses is
differentiable in quadratic mean, so \Cref{thm:cg-fisher-contraction} gives the score projection
and the Fisher loss in \eqref{eq:rg-thermo-finite-missing-information}. For a one-point alphabet
every record row equals one, so $p_\theta(o)=\sum_{\gamma,x}P_\theta(\gamma,x)=1$. The pointwise
observed-information form of this identity is due to \citet{Louis1982}; its average over the record
is \eqref{eq:rg-thermo-finite-missing-information}. $\square$

\noindent\emph{Proof of \Cref{prop:rg-thermo-account-word}.}
The score telescope comes first. Differentiating $P_j=P_{j-1}C_j$ at positive atoms gives
$\mathbf s_j(X_j)=\E[\mathbf s_{j-1}(X_{j-1})\given X_j]$, so the single-step loss is $D_j\succeq0$
by \Cref{thm:cg-fisher-contraction}, and the sum telescopes. The reversed path is Markov, so
$\E[\varepsilon^{\mathrm{sc}}_j\given X_j,\ldots,X_m]=0$; for $j<k$ the increment
$\varepsilon^{\mathrm{sc}}_k$ is measurable with respect to $X_j,\ldots,X_m$, which gives
orthogonality, and $\E[\varepsilon^{\mathrm{sc}}_j(\varepsilon^{\mathrm{sc}}_j)^\top]
=g_{j-1}-g_j$ because $\E[\mathbf s_{j-1}\mathbf s_j^\top]=g_j$. Domination persists under each
kernel on finite spaces, so every divergence is finite, and
\Cref{thm:rg-exact-coarse-vfe,thm:cg-evidence-preserving-channel} at each step give
$\Fenergy_{j-1}-\Fenergy_j=K^{\mathrm{KL}}_j$ at unchanged evidence; summing proves the free-energy
telescope, and division by $N_m$ gives the per-agent form.
$\square$

\noindent\emph{Proof of \Cref{thm:rg-thermo-label-strict}.}
Uniform convergence of $\lambda/N$ gives
$\lambda_{\gamma,N}-\lambda_{\gamma_*,N}\leq-\delta_{\mathsf K}N/2$ for each $\gamma\neq\gamma_*$ and
large $N$; the softmax denominator contains the winning mass, so
$W_{\gamma,N}\leq e^{-\delta_{\mathsf K}N/2}$. Suppress $N$ and write
$\bar\xi_i=\sum_\eta W_\eta\xi_{\eta,i}$. The softmax derivatives
$\partial_iW_\gamma=W_\gamma(\xi_{\gamma,i}-\bar\xi_i)$ and
$\partial_j\bar\xi_i=\sum_\eta W_\eta\partial_j\xi_{\eta,i}+\operatorname{Cov}_W(\xi_i,\xi_j)$ show
that, for $\lvert\beta\rvert\leq2$, $\partial^\beta W_\gamma$ is $W_\gamma$ times a polynomial in
the weights and in the components of the $\xi_\eta$ and $\partial\xi_\eta$, of degree at most
$\lvert\beta\rvert$ in the latter, each of which is $O(N)$ by
\eqref{eq:rg-thermo-label-effective-growth}. Hence
$\|\partial^\beta W_{\gamma,N}\|_{\mathsf K}\leq\mathrm C_{\mathsf K}N^{\lvert\beta\rvert}
e^{-\delta_{\mathsf K}N/2}$ for $\gamma\neq\gamma_*$ and $\lvert\beta\rvert\leq2$, and the same
bound holds for $W_{\gamma_*,N}-1$, which is minus the sum of the losing weights. With
$\delta\xi_\gamma=\xi_\gamma-\xi_{\gamma_*}$ and
$\bar d=\sum_{\gamma\neq\gamma_*}W_\gamma\delta\xi_\gamma$, exact centering gives
$\operatorname{Cov}_W(\xi)=\sum_{\gamma\neq\gamma_*}W_\gamma\delta\xi_\gamma\delta\xi_\gamma^\top-\bar d\bar d^\top$.
Each derivative of $\delta\xi_\gamma$ through order two is $O(N)$, so a term with $j$ derivatives in
the sum is bounded by $\mathrm C_{\mathsf K}N^{j+2}e^{-\delta_{\mathsf K}N/2}$: at most $j$
derivatives fall on a weight and the two score factors contribute at most $N$ each. The product
$\bar d\bar d^\top$ carries two exponential factors and obeys the same bound. Dividing by $N$ proves
\eqref{eq:rg-thermo-label-strict-covariance-bound}. The branch error equals
$\sum_{\gamma\neq\gamma_*}W_\gamma(\mathcal I_{\gamma,N}-\mathcal I_{\gamma_*,N})$, whose branch
differences have bounded $C^2$ norms, so the Leibniz sum with the weight bound gives the stated
bound. Adding the $C^2$ error of the
winning branch and the covariance error, via \eqref{eq:rg-thermo-finite-labeled-fisher}, proves the
$C^2$ convergence, since exponentials dominate the powers of $N$. For the floor, uniform convergence
in operator norm gives $\|g_N-\mathcal I_{\gamma_*}\|\leq\tfrac12\mathrm c_{\min}$ on $\mathsf K$ for
large $N$, so $z^\top g_Nz\geq z^\top\mathcal I_{\gamma_*}z-\tfrac12\mathrm c_{\min}\lvert z\rvert^2
\geq\tfrac12\mathrm c_{\min}\lvert z\rvert^2$. $\square$

\noindent\emph{Proof of \Cref{prop:rg-thermo-label-tie-scale}.}
The label set is finite, so there is a positive bulk gap from $S$ to every other label at
$\theta_0$; comparison with one maximizing label and convergence of $\lambda/N$ make every
nonmaximizing weight exponentially small, so its limit is zero. $C^1$ convergence gives
$\xi_{\gamma,N}/N\to\mathsf v_\gamma$. Dividing \eqref{eq:rg-thermo-finite-labeled-fisher} by $N^2$
makes the branch term $O(1/N)$ and leaves $\operatorname{Cov}_{W_N}(\xi_{\gamma,N}/N)$, which
converges by finite summation. For every $z$,
$z^\top\operatorname{Cov}_{\bar{\mathsf p}}(\mathsf v_\gamma)z
=\sum_\gamma\bar{\mathsf p}_\gamma[z\cdot(\mathsf v_\gamma-\bar{\mathsf v})]^2$, which vanishes
exactly when $z\cdot\mathsf v_\gamma$ is constant on the positive-weight labels, so the kernel is the
orthogonal complement of the stated span. A positive limiting quadratic form gives the asymptotic
after multiplication by $N^2$, and a zero one gives only the little-o conclusion. $\square$

\noindent\emph{Proof of \Cref{thm:rg-thermo-label-crossover}.}
A compact set of $\mathsf w$ is bounded, so its image under $\theta_N$ and the segments to
$\theta_0$ lie in the chosen neighborhood for large $N$. Put
$\lambda^{\mathrm{inc}}_{\gamma,N}(\mathsf w)=\lambda_{\gamma,N}(\theta_0+\mathsf w/N)-\lambda_{\gamma,N}(\theta_0)$.
The fundamental theorem of calculus gives
\begin{equation}
\lambda^{\mathrm{inc}}_{\gamma,N}(\mathsf w)
=\int_0^1N^{-1}\nabla\lambda_{\gamma,N}(\theta_0+\tau\mathsf w/N)\cdot\mathsf w d\tau,
\label{eq:rg-thermo-label-window-increment}
\end{equation}
whose integrand tends uniformly to $\mathsf v_\gamma\cdot\mathsf w$ by the $C^1$ part of
\eqref{eq:rg-thermo-label-effective-growth} and continuity of $\nabla\mathsf s_\gamma$.
Differentiation in $\mathsf w$ gives
\begin{equation}
D_{\mathsf w}\lambda^{\mathrm{inc}}_{\gamma,N}\to\mathsf v_\gamma,
\qquad
\|D^2_{\mathsf w}\lambda^{\mathrm{inc}}_{\gamma,N}\|=O(N^{-1}),
\qquad
\|D^3_{\mathsf w}\lambda^{\mathrm{inc}}_{\gamma,N}\|=O(N^{-2}),
\label{eq:rg-thermo-label-window-derivatives}
\end{equation}
so $\lambda^{\mathrm{inc}}_{\gamma,N}\to\mathsf w\cdot\mathsf v_\gamma$ in $C^3$. For tied labels,
the relative log masses $\lambda_{\gamma,N}-\lambda_{\gamma_0,N}$ at $\theta_N(\mathsf w)$ are the
logarithms of the center ratios plus the difference of increments, and converge in $C^3$ to
$\log\mathrm c_\gamma+\mathsf w\cdot(\mathsf v_\gamma-\mathsf v_{\gamma_0})$; exponentiation and the
finite softmax are smooth on the resulting compact ranges, with denominator at least the reference
mass one. For a nonmaximizing label the relative center log mass is at most $-\mathrm cN$ for some
constant $\mathrm c>0$, while its increment and $\mathsf w$-derivatives through order three are
bounded, so its relative exponential and those derivatives are $O(e^{-\mathrm cN})$. Normalizing
proves \eqref{eq:rg-thermo-label-limiting-weights} in $C^3$, including the zero coordinates outside
$S$. With $\hat\xi_{\gamma,N}(\mathsf w)=N^{-1}\xi_{\gamma,N}(\theta_N(\mathsf w))$,
\eqref{eq:rg-thermo-label-window-derivatives} gives $\hat\xi_{\gamma,N}\to\mathsf v_\gamma$ in $C^2$,
and the affine change $\theta_N$ has derivative $N^{-1}\mathrm{Id}$, so
\eqref{eq:rg-thermo-finite-labeled-fisher} becomes
\begin{equation}
g^F_{\mathsf w}(\mathcal E_N)(\mathsf w)
=\operatorname{Cov}_{W_N(\theta_N(\mathsf w))}\bigl(\hat\xi_{\gamma,N}(\mathsf w)\bigr)
+\frac1N\sum_\gamma W_{\gamma,N}(\theta_N(\mathsf w))\mathcal I_{\gamma,N}(\theta_N(\mathsf w)).
\label{eq:rg-thermo-label-window-decomposition}
\end{equation}
The covariance converges in $C^2$ by finite sums and products, without dividing by a weight that
tends to zero. The branch densities and their first two $\mathsf w$-derivatives are bounded, since a
$\mathsf w$-derivative is $N^{-1}$ times a $\theta$-derivative, and the weights have bounded $C^2$
norms, so every term of the Leibniz sum for the last term is $O(N^{-1})$. Finally the
log-sum-exponential has gradient $\sum_\gamma\bar{\mathsf p}_\gamma\mathsf v_\gamma$ and Hessian
$\operatorname{Cov}_{\bar{\mathsf p}}(\mathsf v_\gamma)$. The window is the finite-mixture
counterpart of the finite-volume rounding of a first-order transition, which
\citet{BorgsKotecky1990} derive from a lattice model as a hyperbolic-tangent crossover of width of
order the inverse volume. $\square$

\noindent\emph{Proof of \Cref{prop:rg-thermo-label-tie-counterexamples}.}
The common branch law cancels from the relative log mass and has zero score, so the label is the
only source of information. Then $W_1'=W_1(1-W_1)\omega'_N$, the centered label scores are
$(1-W_1)\omega'_N$ and $-W_1\omega'_N$, and averaging their squares gives
\eqref{eq:rg-thermo-label-binary-formula}, with $W_1(1-W_1)=\operatorname{bv}(\omega_N)$. In (a) the
$C^3$ norm of $\omega_N$ on a compact interval is at most $\mathrm CN$ and $\omega_N/N=\theta$; the
formulas follow by substitution. In (b) the derivatives vanish and $\omega_N/N\to0$ in $C^3$. In (c)
the $C^3$ norms are $O(N^{\mathrm c})$ and $O(\sqrt N)$, both $o(N)$, and substitution gives the two
statements, with $\operatorname{bv}(0)=1/4$ and $\operatorname{bv}(\log2)=2/9$. In (d)
$\omega_N/N=\theta-\mathrm cN^{-1}\log N\to\theta$ in $C^3$ and the $C^3$ norm is $O(N)$; at zero,
$\operatorname{bv}(-\mathrm c\log N)N^2/N=N^{1+\mathrm c}/(1+N^{\mathrm c})^2$. In (e) derivatives
through order three are bounded by $N^{-1/6}N=N^{5/6}$, so $\omega_N/N\to0$ in $C^3$;
$\lvert\omega_N\rvert\leq N^{-1/6}$ gives uniform convergence of the weights, and
$0\leq F_N\leq N^{-1}\tfrac14N^{1/3}$. The Leibniz rule gives
\eqref{eq:rg-thermo-label-oscillatory-derivative}. At zero $\omega_N=0$, $\omega'_N=N^{1/6}$,
$\omega''_N=0$, $\omega'''_N=-N^{5/6}$, and $\operatorname{bv}(0)=1/4$, $\operatorname{bv}'(0)=0$,
$\operatorname{bv}''(0)=-1/8$, which gives the value of $F''_N(0)$; the uniform limit has second
derivative zero. $\square$
